% Requires \usepackage{graphicx,amsmath}.
\begin{figure*}[t]
\centering
\includegraphics[width=\textwidth]{Figure_B_historical_QF.pdf}
\caption{Historical ImageNet-A seed0 Task0 neighborhood measurements for EWCLoRA (SGD) and FS-EWCLoRA (visual GAM). From left: directional mean positive Task0 training CE damage, directional mean training accuracy (\%), and directional mean fixed-anchor trace-normalized Fisher quadratic. $\overline Q_F(r)=K^{-1}\sum_k\frac12r^2u_k^\top\widetilde F_0u_k$, where $K=100$ and $\widetilde F_0$ is the stored diagonal anchor Fisher proxy with trace approximately one. This proxy squares minibatch-mean CE gradients, is fixed throughout the radius scan and has no EWC lambda multiplier; it is not raw per-example empirical Fisher trace at the perturbed point. All 81 radii from 0 to 20 in steps of 0.25 are plotted; markers are shown every 2.5 units. Fresh isotropic directions were drawn at each radius; cross-method direction pairing is unverified. All three quantities are read from the same historical training-split radius/direction records; train inputs use one cached draw of stochastic training augmentation, and each method retains its own fixed Task0 head. Positive damage is $[L_0(\theta_0+ru)-L_0(\theta_0)]_+$. Weight drift is the joint Frobenius radius over 24 effective Q/V matrices (14,155,776 coordinates). The nearly overlapping third-panel means are expected for isotropic directions because the anchor traces were normalized. This historical plot uses a different measurement protocol from the deterministic raw-Fisher diagnostic; its curves are not mixed with those new measurements.}
\end{figure*}
